% --- Archon physics formalization source begin ---
% archon:physics
% archon:covers PhyXMiniProblems/problem_phyx_mini_0229.lean
% archon:source-report reports/phyx_mini/problem_phyx_mini_0229.source.json

\chapter{Physics problem phyx\_mini\_0229}
\label{ch:PhyXMiniProblems_problem_phyx_mini_0229}

\paragraph{Problem source.}
\#\# Physical scenario

A small ball of mass \textbackslash{}( M \textbackslash{}) is attached to the end of a uniform rod of equal mass \textbackslash{}( M \textbackslash{}) and length \textbackslash{}( L \textbackslash{}) that is pivoted at the top (figure). \textbackslash{}( L = 2.00\textbackslash{},\textbackslash{}mathrm\{m\} \textbackslash{}).

\#\# Question

Determine this period.

\#\# Answer choices

- A: A: 5.68s

- B: B: 4.68s

- C: C: 3.68s

- D: D: 2.68s

\#\# Figure caption (auxiliary; use the image as primary evidence)

The image depicts a vertical rod with a pivot at the top. The rod has a point labeled \textbackslash{}( P \textbackslash{}) marked on it. The distance from the pivot to point \textbackslash{}( P \textbackslash{}) is indicated as \textbackslash{}( y \textbackslash{}), and the distance from the pivot to the bottom of the rod is labeled \textbackslash{}( L \textbackslash{}). At the bottom of the rod, there is a mass labeled \textbackslash{}( M \textbackslash{}). The horizontal line at the bottom is marked as \textbackslash{}( y = 0 \textbackslash{}). There are arrows indicating the measurements \textbackslash{}( L \textbackslash{}) and \textbackslash{}( y \textbackslash{}).

\#\# Dataset metadata

Resonance and Harmonics; Physical Model Grounding Reasoning, Numerical Reasoning

\paragraph{Recorded answer/context.}
D: D: 2.68s

\paragraph{Figure/image path.}
/root/proposal\_for\_physic/science-mango/phyx\_mini\_run/phyx\_data/test\_image/229.png

\paragraph{Formalization target.}
create a compiling Lean file with sorry bodies at `PhyXMiniProblems/problem\_phyx\_mini\_0229.lean`. The Lean declarations must preserve the physical quantities, dimensions or dimensional roles, figure labels, governing-law hypotheses, and final relation expressed by this problem.
Use Mathlib/Physlib names found through LeanExplore where available. If a physics API is missing, introduce faithful local abstractions rather than scalar placeholder aliases.

\begin{theorem}[Physics formalization target]
\label{thm:physics:phyx_mini_0229:target}
\lean{PhyXMiniProblems.ProblemPhyXMini0229.rodAndEqualBallLinearizedPeriod_matches_recordedAnswerD}
\uses{def:physics:phyx-mini-0229:phyxminiproblems-problemphyxmini0229-lengthinmeters, def:physics:phyx-mini-0229:phyxminiproblems-problemphyxmini0229-timeinseconds, def:physics:phyx-mini-0229:phyxminiproblems-problemphyxmini0229-accelerationinmeterspersecondsquared, def:physics:phyx-mini-0229:phyxminiproblems-problemphyxmini0229-rodandballpendulumsetup, def:physics:phyx-mini-0229:phyxminiproblems-problemphyxmini0229-matchesprimarypendulumfigure, def:physics:phyx-mini-0229:phyxminiproblems-problemphyxmini0229-hasrodandballpendulumproblemdata, def:physics:phyx-mini-0229:phyxminiproblems-problemphyxmini0229-satisfiesrodandballcompoundpendulumlaws, def:physics:phyx-mini-0229:phyxminiproblems-problemphyxmini0229-answerchoice, def:physics:phyx-mini-0229:phyxminiproblems-problemphyxmini0229-recordedanswerchoice, def:physics:phyx-mini-0229:phyxminiproblems-problemphyxmini0229-matchesdisplayedperiod}
The assigned autoformalize agent should translate this physics problem into Lean declarations in the covered file, with theorem and lemma proof bodies written as `by sorry`.
\end{theorem}
\begin{proof}
This is an autoformalization task, not a proof task. Produce statements that can later be proved by the physics prover without weakening the physical model.
\end{proof}

% --- Archon declaration topology begin ---
\section{Declaration topology}
\begin{definition}[Mass Quantity]
  \label{def:physics:phyx-mini-0229:phyxminiproblems-problemphyxmini0229-massquantity}
  \lean{PhyXMiniProblems.ProblemPhyXMini0229.MassQuantity}
  A nonnegative physical mass
\end{definition}
\begin{proof}
  This is the declared model definition of Mass Quantity.
\end{proof}
\begin{definition}[Length Quantity]
  \label{def:physics:phyx-mini-0229:phyxminiproblems-problemphyxmini0229-lengthquantity}
  \lean{PhyXMiniProblems.ProblemPhyXMini0229.LengthQuantity}
  A nonnegative physical length
\end{definition}
\begin{proof}
  This is the declared model definition of Length Quantity.
\end{proof}
\begin{definition}[Time Quantity]
  \label{def:physics:phyx-mini-0229:phyxminiproblems-problemphyxmini0229-timequantity}
  \lean{PhyXMiniProblems.ProblemPhyXMini0229.TimeQuantity}
  A nonnegative physical duration
\end{definition}
\begin{proof}
  This is the declared model definition of Time Quantity.
\end{proof}
\begin{definition}[Acceleration Magnitude Quantity]
  \label{def:physics:phyx-mini-0229:phyxminiproblems-problemphyxmini0229-accelerationmagnitudequantity}
  \lean{PhyXMiniProblems.ProblemPhyXMini0229.AccelerationMagnitudeQuantity}
  A nonnegative acceleration magnitude, with dimension L T⁻²
\end{definition}
\begin{proof}
  This is the declared model definition of Acceleration Magnitude Quantity.
\end{proof}
\begin{definition}[Moment Of Inertia Quantity]
  \label{def:physics:phyx-mini-0229:phyxminiproblems-problemphyxmini0229-momentofinertiaquantity}
  \lean{PhyXMiniProblems.ProblemPhyXMini0229.MomentOfInertiaQuantity}
  A moment of inertia about the pivot, with dimension M L²
\end{definition}
\begin{proof}
  This is the declared model definition of Moment Of Inertia Quantity.
\end{proof}
\begin{definition}[Restoring Torque Coefficient Quantity]
  \label{def:physics:phyx-mini-0229:phyxminiproblems-problemphyxmini0229-restoringtorquecoefficientquantity}
  \lean{PhyXMiniProblems.ProblemPhyXMini0229.RestoringTorqueCoefficientQuantity}
  The nonnegative coefficient κ in the exact restoring torque τ(θ) = -κ sin θ. Since radians are dimensionless, κ has the torque dimension M L² T⁻²
\end{definition}
\begin{proof}
  This is the declared model definition of Restoring Torque Coefficient Quantity.
\end{proof}
\begin{definition}[Torque Quantity]
  \label{def:physics:phyx-mini-0229:phyxminiproblems-problemphyxmini0229-torquequantity}
  \lean{PhyXMiniProblems.ProblemPhyXMini0229.TorqueQuantity}
  A signed torque, with dimension M L² T⁻²
\end{definition}
\begin{proof}
  This is the declared model definition of Torque Quantity.
\end{proof}
\begin{definition}[mass In Kilograms]
  \label{def:physics:phyx-mini-0229:phyxminiproblems-problemphyxmini0229-massinkilograms}
  \lean{PhyXMiniProblems.ProblemPhyXMini0229.massInKilograms}
  \uses{def:physics:phyx-mini-0229:phyxminiproblems-problemphyxmini0229-massquantity}
  Read a physical mass in kilograms
\end{definition}
\begin{proof}
  This is the declared model definition of mass In Kilograms.
\end{proof}
\begin{definition}[length In Meters]
  \label{def:physics:phyx-mini-0229:phyxminiproblems-problemphyxmini0229-lengthinmeters}
  \lean{PhyXMiniProblems.ProblemPhyXMini0229.lengthInMeters}
  \uses{def:physics:phyx-mini-0229:phyxminiproblems-problemphyxmini0229-lengthquantity}
  Read a physical length in meters
\end{definition}
\begin{proof}
  This is the declared model definition of length In Meters.
\end{proof}
\begin{definition}[time In Seconds]
  \label{def:physics:phyx-mini-0229:phyxminiproblems-problemphyxmini0229-timeinseconds}
  \lean{PhyXMiniProblems.ProblemPhyXMini0229.timeInSeconds}
  \uses{def:physics:phyx-mini-0229:phyxminiproblems-problemphyxmini0229-timequantity}
  Read a physical duration in seconds
\end{definition}
\begin{proof}
  This is the declared model definition of time In Seconds.
\end{proof}
\begin{definition}[acceleration In Meters Per Second Squared]
  \label{def:physics:phyx-mini-0229:phyxminiproblems-problemphyxmini0229-accelerationinmeterspersecondsquared}
  \lean{PhyXMiniProblems.ProblemPhyXMini0229.accelerationInMetersPerSecondSquared}
  \uses{def:physics:phyx-mini-0229:phyxminiproblems-problemphyxmini0229-accelerationmagnitudequantity}
  Read an acceleration magnitude in meters per second squared
\end{definition}
\begin{proof}
  This is the declared model definition of acceleration In Meters Per Second Squared.
\end{proof}
\begin{definition}[moment Of Inertia In Kilogram Meters Squared]
  \label{def:physics:phyx-mini-0229:phyxminiproblems-problemphyxmini0229-momentofinertiainkilogrammeterssquared}
  \lean{PhyXMiniProblems.ProblemPhyXMini0229.momentOfInertiaInKilogramMetersSquared}
  \uses{def:physics:phyx-mini-0229:phyxminiproblems-problemphyxmini0229-momentofinertiaquantity}
  Read a moment of inertia in kilogram meter squared
\end{definition}
\begin{proof}
  This is the declared model definition of moment Of Inertia In Kilogram Meters Squared.
\end{proof}
\begin{definition}[restoring Coefficient In Newton Meters]
  \label{def:physics:phyx-mini-0229:phyxminiproblems-problemphyxmini0229-restoringcoefficientinnewtonmeters}
  \lean{PhyXMiniProblems.ProblemPhyXMini0229.restoringCoefficientInNewtonMeters}
  \uses{def:physics:phyx-mini-0229:phyxminiproblems-problemphyxmini0229-restoringtorquecoefficientquantity}
  Read a restoring-torque coefficient in newton meters
\end{definition}
\begin{proof}
  This is the declared model definition of restoring Coefficient In Newton Meters.
\end{proof}
\begin{definition}[torque In Newton Meters]
  \label{def:physics:phyx-mini-0229:phyxminiproblems-problemphyxmini0229-torqueinnewtonmeters}
  \lean{PhyXMiniProblems.ProblemPhyXMini0229.torqueInNewtonMeters}
  \uses{def:physics:phyx-mini-0229:phyxminiproblems-problemphyxmini0229-torquequantity}
  Read a signed torque in newton meters
\end{definition}
\begin{proof}
  This is the declared model definition of torque In Newton Meters.
\end{proof}
\begin{definition}[Pendulum Component]
  \label{def:physics:phyx-mini-0229:phyxminiproblems-problemphyxmini0229-pendulumcomponent}
  \lean{PhyXMiniProblems.ProblemPhyXMini0229.PendulumComponent}
  The two massive components of the compound pendulum
\end{definition}
\begin{proof}
  This is the declared model definition of Pendulum Component.
\end{proof}
\begin{definition}[Figure Point]
  \label{def:physics:phyx-mini-0229:phyxminiproblems-problemphyxmini0229-figurepoint}
  \lean{PhyXMiniProblems.ProblemPhyXMini0229.FigurePoint}
  Distinguished points shown in the primary figure
\end{definition}
\begin{proof}
  This is the declared model definition of Figure Point.
\end{proof}
\begin{definition}[Figure Label]
  \label{def:physics:phyx-mini-0229:phyxminiproblems-problemphyxmini0229-figurelabel}
  \lean{PhyXMiniProblems.ProblemPhyXMini0229.FigureLabel}
  Text or mathematical labels visible in the primary figure
\end{definition}
\begin{proof}
  This is the declared model definition of Figure Label.
\end{proof}
\begin{definition}[Rod Geometry]
  \label{def:physics:phyx-mini-0229:phyxminiproblems-problemphyxmini0229-rodgeometry}
  \lean{PhyXMiniProblems.ProblemPhyXMini0229.RodGeometry}
  The idealized geometry of the extended component
\end{definition}
\begin{proof}
  This is the declared model definition of Rod Geometry.
\end{proof}
\begin{definition}[End Ball Idealization]
  \label{def:physics:phyx-mini-0229:phyxminiproblems-problemphyxmini0229-endballidealization}
  \lean{PhyXMiniProblems.ProblemPhyXMini0229.EndBallIdealization}
  The small attached ball's dynamical idealization
\end{definition}
\begin{proof}
  This is the declared model definition of End Ball Idealization.
\end{proof}
\begin{definition}[Rod And Ball Pendulum Setup]
  \label{def:physics:phyx-mini-0229:phyxminiproblems-problemphyxmini0229-rodandballpendulumsetup}
  \lean{PhyXMiniProblems.ProblemPhyXMini0229.RodAndBallPendulumSetup}
  \uses{def:physics:phyx-mini-0229:phyxminiproblems-problemphyxmini0229-massquantity, def:physics:phyx-mini-0229:phyxminiproblems-problemphyxmini0229-lengthquantity, def:physics:phyx-mini-0229:phyxminiproblems-problemphyxmini0229-timequantity, def:physics:phyx-mini-0229:phyxminiproblems-problemphyxmini0229-accelerationmagnitudequantity, def:physics:phyx-mini-0229:phyxminiproblems-problemphyxmini0229-momentofinertiaquantity, def:physics:phyx-mini-0229:phyxminiproblems-problemphyxmini0229-restoringtorquecoefficientquantity, def:physics:phyx-mini-0229:phyxminiproblems-problemphyxmini0229-torquequantity, def:physics:phyx-mini-0229:phyxminiproblems-problemphyxmini0229-pendulumcomponent, def:physics:phyx-mini-0229:phyxminiproblems-problemphyxmini0229-figurepoint, def:physics:phyx-mini-0229:phyxminiproblems-problemphyxmini0229-figurelabel, def:physics:phyx-mini-0229:phyxminiproblems-problemphyxmini0229-rodgeometry, def:physics:phyx-mini-0229:phyxminiproblems-problemphyxmini0229-endballidealization}
  Physical quantities and labeled geometry for the rod--ball compound pendulum. The coordinate verticalHeightAboveYZero follows the primary image: the lower-end horizontal line is y = 0, while the marked point P has height y. The separate distanceBelowPivot supplies lever arms. The nonlinear period function is constrained only locally near zero peak angle by the laws below; no arbitrary-amplitude small-angle equality is built into this setup
\end{definition}
\begin{proof}
  This is the declared model definition of Rod And Ball Pendulum Setup.
\end{proof}
\begin{definition}[Matches Primary Pendulum Figure]
  \label{def:physics:phyx-mini-0229:phyxminiproblems-problemphyxmini0229-matchesprimarypendulumfigure}
  \lean{PhyXMiniProblems.ProblemPhyXMini0229.MatchesPrimaryPendulumFigure}
  \uses{def:physics:phyx-mini-0229:phyxminiproblems-problemphyxmini0229-lengthinmeters, def:physics:phyx-mini-0229:phyxminiproblems-problemphyxmini0229-rodandballpendulumsetup}
  Primary-image geometry and labels. In particular, the pivot is one rod length above the y = 0 line, the ball is attached at the lower end, and P lies strictly inside the rod. The complementary-coordinate identity records the image's actual convention: y is measured upward from the lower end
\end{definition}
\begin{proof}
  This is the declared model definition of Matches Primary Pendulum Figure.
\end{proof}
\begin{definition}[Has Rod And Ball Pendulum Problem Data]
  \label{def:physics:phyx-mini-0229:phyxminiproblems-problemphyxmini0229-hasrodandballpendulumproblemdata}
  \lean{PhyXMiniProblems.ProblemPhyXMini0229.HasRodAndBallPendulumProblemData}
  \uses{def:physics:phyx-mini-0229:phyxminiproblems-problemphyxmini0229-massinkilograms, def:physics:phyx-mini-0229:phyxminiproblems-problemphyxmini0229-lengthinmeters, def:physics:phyx-mini-0229:phyxminiproblems-problemphyxmini0229-accelerationinmeterspersecondsquared, def:physics:phyx-mini-0229:phyxminiproblems-problemphyxmini0229-rodandballpendulumsetup}
  Problem and numerical-evaluation data. The rod and ball have the same positive mass, the length is 2.00 m, and the standard textbook value g = 9.8 m/s² is the environmental readout used to select among the choices. The slender-rod and point-ball fields state the mechanical idealizations implicit in the problem's wording
\end{definition}
\begin{proof}
  This is the declared model definition of Has Rod And Ball Pendulum Problem Data.
\end{proof}
\begin{definition}[Satisfies Rod And Ball Compound Pendulum Laws]
  \label{def:physics:phyx-mini-0229:phyxminiproblems-problemphyxmini0229-satisfiesrodandballcompoundpendulumlaws}
  \lean{PhyXMiniProblems.ProblemPhyXMini0229.SatisfiesRodAndBallCompoundPendulumLaws}
  \uses{def:physics:phyx-mini-0229:phyxminiproblems-problemphyxmini0229-massinkilograms, def:physics:phyx-mini-0229:phyxminiproblems-problemphyxmini0229-lengthinmeters, def:physics:phyx-mini-0229:phyxminiproblems-problemphyxmini0229-timeinseconds, def:physics:phyx-mini-0229:phyxminiproblems-problemphyxmini0229-accelerationinmeterspersecondsquared, def:physics:phyx-mini-0229:phyxminiproblems-problemphyxmini0229-momentofinertiainkilogrammeterssquared, def:physics:phyx-mini-0229:phyxminiproblems-problemphyxmini0229-restoringcoefficientinnewtonmeters, def:physics:phyx-mini-0229:phyxminiproblems-problemphyxmini0229-torqueinnewtonmeters, def:physics:phyx-mini-0229:phyxminiproblems-problemphyxmini0229-rodandballpendulumsetup}
  Moment-of-inertia, gravitational-restoring, and local oscillation laws for the compound pendulum. The first five relations are the rigid-body and point-mass laws. The exact signed torque is -κ sin θ. Its derivative and little-o remainder are stated at θ = 0, so replacing sin θ by θ is explicitly local. The angular frequency and period laws concern the resulting linearized oscillator.
\end{definition}
\begin{proof}
  This is the declared model definition of Satisfies Rod And Ball Compound Pendulum Laws.
\end{proof}
\begin{definition}[Answer Choice]
  \label{def:physics:phyx-mini-0229:phyxminiproblems-problemphyxmini0229-answerchoice}
  \lean{PhyXMiniProblems.ProblemPhyXMini0229.AnswerChoice}
  Labels of the four displayed period choices
\end{definition}
\begin{proof}
  This is the declared model definition of Answer Choice.
\end{proof}
\begin{definition}[seconds]
  \label{def:physics:phyx-mini-0229:phyxminiproblems-problemphyxmini0229-answerchoice-seconds}
  \lean{PhyXMiniProblems.ProblemPhyXMini0229.AnswerChoice.seconds}
  \uses{def:physics:phyx-mini-0229:phyxminiproblems-problemphyxmini0229-answerchoice}
  The period in seconds printed beside each answer label
\end{definition}
\begin{proof}
  This is the declared model definition of seconds.
\end{proof}
\begin{definition}[recorded Answer Choice]
  \label{def:physics:phyx-mini-0229:phyxminiproblems-problemphyxmini0229-recordedanswerchoice}
  \lean{PhyXMiniProblems.ProblemPhyXMini0229.recordedAnswerChoice}
  \uses{def:physics:phyx-mini-0229:phyxminiproblems-problemphyxmini0229-answerchoice}
  The answer label recorded in the supplied dataset
\end{definition}
\begin{proof}
  This is the declared model definition of recorded Answer Choice.
\end{proof}
\begin{definition}[Matches Displayed Period]
  \label{def:physics:phyx-mini-0229:phyxminiproblems-problemphyxmini0229-matchesdisplayedperiod}
  \lean{PhyXMiniProblems.ProblemPhyXMini0229.MatchesDisplayedPeriod}
  \uses{def:physics:phyx-mini-0229:phyxminiproblems-problemphyxmini0229-timequantity, def:physics:phyx-mini-0229:phyxminiproblems-problemphyxmini0229-timeinseconds, def:physics:phyx-mini-0229:phyxminiproblems-problemphyxmini0229-answerchoice}
  Agreement with a displayed value after rounding to two decimal places
\end{definition}
\begin{proof}
  This is the declared model definition of Matches Displayed Period.
\end{proof}
% --- Archon declaration topology end ---
% --- Archon physics formalization source end ---
